% ====================================================================
%  ISI MODUL BAB 2: KOMBINATORIKA
%  Jangan dikompilasi langsung. Kompilasi salah satu dari:
%    kombin_siswa.tex  -> Edisi Siswa
%    kombin_guru.tex   -> Edisi Guru (ada catatan mengajar & pembahasan)
% ====================================================================
\makeatletter
\@ifundefined{ifguru}{\newif\ifguru}{}
\makeatother

\documentclass[11pt,a4paper]{article}
% ====================================================================
%  GAYA BERSAMA MODUL OLIMPIADE (dipakai teobil_isi.tex & kombin_isi.tex)
%  Semua tampilan hitam-putih / abu-abu supaya hemat tinta.
% ====================================================================
\usepackage[bahasai]{babel}
\usepackage{amsmath, amssymb}
\usepackage{geometry}
\geometry{top=2.5cm, bottom=2.5cm, left=2.5cm, right=2.5cm}
\usepackage[most]{tcolorbox}
\usepackage{enumitem}
\usepackage{array}
\usepackage{comment}
\usepackage{tikz}

\setlength{\parindent}{0pt}
\setlength{\parskip}{5pt}

\tcbset{colback=white, colframe=black, colbacktitle=white, coltitle=black,
  fonttitle=\bfseries, boxrule=0.6pt, arc=1mm,
  before upper={\setlength{\parskip}{5pt}}}

% Kegiatan pembuka
\newtcolorbox{selidiki}{enhanced, breakable, colback=black!4, colframe=black!60,
  colbacktitle=black!15, coltitle=black, title={Ayo Selidiki!}}

% Contoh soal (abu-abu tipis, garis tebal di kiri)
\newtcolorbox{contoh}{enhanced, breakable, colback=black!4, colframe=black!55,
  boxrule=0pt, leftrule=3pt, arc=0pt, left=3mm, before skip=6pt, after skip=6pt}

% Soal untuk dicoba langsung (putih, garis putus-putus)
\newtcolorbox{giliranmu}{enhanced, borderline={0.6pt}{0pt}{black, dashed},
  frame hidden, title={Sekarang Giliranmu!},
  attach boxed title to top left={xshift=3mm, yshift=-2mm},
  boxed title style={colframe=white, colback=white}}

% Ringkasan akhir topik
\newtcolorbox{rangkuman}{enhanced, colback=white, colframe=black!60,
  colbacktitle=black!15, coltitle=black, title={Rangkuman}}

% Garis isian
\newcommand{\kosong}{\rule{1.2cm}{0.4pt}}

% Label level: blok abu-abu kecil
\newcommand{\level}[1]{\par\medskip\noindent\colorbox{black!12}{\strut\textbf{#1}}\par\smallskip}

% Pembahasan soal: hanya muncul di Edisi Guru
\newtcolorbox{kotakjawab}{enhanced, breakable, boxrule=0pt, leftrule=2pt, arc=0pt,
  colback=white, colframe=black, left=2mm, top=1mm, bottom=1mm,
  before skip=3pt, after skip=3pt}
\newcommand{\jawab}[2]{\ifguru
  \begin{kotakjawab}\small\textbf{Jawaban: #1.} #2\end{kotakjawab}\fi}

% Catatan mengajar: hanya muncul di Edisi Guru
\ifguru
  \newtcolorbox{catatanguru}{breakable, colback=white, colframe=black!60,
    colbacktitle=black!15, title={Catatan untuk Guru}}
\else
  \excludecomment{catatanguru}
\fi

\newcommand{\C}[2]{\binom{#1}{#2}}


\title{\textbf{Modul Pembinaan Olimpiade Matematika SMP\\Bab 2: Kombinatorika}}
\author{\ifguru Edisi Guru\else Bahan Ajar Siswa (Kelas 7, 8, 9)\fi}
\date{}

\begin{document}

\maketitle
\thispagestyle{empty}

\begin{tcolorbox}[title={Halo lagi, Calon Sang Juara!}]
Di Bab 1 kamu sudah belajar mencari pola pada bilangan. Sekarang kita belajar
\textbf{menghitung banyaknya cara}: berapa banyak susunan, pilihan, atau kemungkinan, tanpa
harus menuliskannya satu per satu. Kalau bingung, jangan langsung cari rumus. Coba dulu dengan
angka kecil, tulis semua kemungkinannya, lalu lihat polanya.
\end{tcolorbox}

\textbf{Cara memakai modul ini.} Sama seperti Bab 1: setiap topik dimulai dengan
\textbf{Ayo Selidiki}, lalu dibagi menjadi beberapa langkah kecil. Di setiap langkah ada
contoh, lalu kamu langsung mencoba sendiri di kotak \textbf{Sekarang Giliranmu}. Di akhir topik
ada latihan soal dari Level 1 sampai Level 5.

\ifguru
\begin{catatanguru}
\textbf{Rencana pertemuan (90 menit per pertemuan):} satu topik untuk satu pertemuan. Topik 4
(Kombinasi) boleh dua pertemuan. Satu pertemuan terakhir untuk Soal Campuran.

\textbf{Saran umum.} Tunda memberi rumus. Biarkan siswa mendaftar semua kemungkinan untuk kasus
kecil di Ayo Selidiki. Rumus baru diberikan setelah mereka merasa ``capek mendaftar''.

\textbf{Level 4 dan 5} adalah soal uraian. Nilailah dari kelengkapan alasan, bukan hanya
jawaban akhir.
\end{catatanguru}
\fi

\newpage
% ====================================================================
\section*{Topik 1: Aturan Penjumlahan dan Perkalian}
% ====================================================================

\begin{catatanguru}
Kesalahan yang paling sering muncul: menjumlah padahal harus mengali, lupa bahwa angka 0 tidak
boleh di depan, dan mengisi kotak dari kiri padahal syaratnya ada di kotak kanan.
\end{catatanguru}

\begin{selidiki}
Di kantin ada 3 makanan (Nasi goreng, Mi ayam, Bakso) dan 4 minuman (Es teh, Es jeruk, Susu,
Air). Kamu membeli satu makanan \textbf{dan} satu minuman. Isi setiap kotak dengan pasangannya,
misalnya ``Bakso + Susu''.

\begin{center}
\renewcommand{\arraystretch}{1.8}
\begin{tabular}{|l|*{4}{>{\centering\arraybackslash}p{2.2cm}|}}
\hline
 & Es teh & Es jeruk & Susu & Air\\ \hline
Nasi goreng & NG + Es teh & & & \\ \hline
Mi ayam & & & & \\ \hline
Bakso & & & & \\ \hline
\end{tabular}
\end{center}

Sekarang jawab:
\begin{enumerate}
    \item Ada berapa pasangan seluruhnya? Bisakah kamu menghitungnya tanpa mengisi tabel?
    \item Kalau kantin menambah 1 minuman, pasangannya bertambah berapa?
    \item Kalau kamu hanya boleh membeli \textbf{satu saja}, makanan atau minuman, ada berapa pilihan?
\end{enumerate}

Tabelnya punya 3 baris dan 4 kolom, jadi ada $3 \times 4 = 12$ pasangan. Tapi kalau hanya
membeli satu saja, pilihannya $3 + 4 = 7$. Kapan dikali dan kapan dijumlah? Itulah yang akan
kita pelajari.
\end{selidiki}

\subsection*{Langkah 1: Aturan perkalian (kata kunci: DAN)}
Kalau sebuah kegiatan dilakukan dalam beberapa tahap, \textbf{kalikan} banyak pilihan di setiap
tahap.

\begin{contoh}
\textbf{Contoh.} Rudi punya 3 baju, 2 celana, dan 2 pasang sepatu. Ada berapa cara ia berpakaian?

\textit{Jawab:} Pilih baju \textbf{dan} celana \textbf{dan} sepatu: $3 \times 2 \times 2 = \textbf{12}$ cara.
\end{contoh}

\begin{giliranmu}
\begin{enumerate}[label=\alph*.]
    \item Toko roti punya 4 rasa roti dan 3 jenis selai. Banyak pasangan roti dan selai \kosong.
    \item Dua koin dan satu dadu dilempar bersamaan. Banyak hasil yang mungkin \kosong.
\end{enumerate}
\end{giliranmu}

\subsection*{Langkah 2: Aturan penjumlahan (kata kunci: ATAU)}
Kalau kita hanya memilih \textbf{salah satu} dari beberapa kelompok, \textbf{jumlahkan} banyak
pilihannya.

\begin{contoh}
\textbf{Contoh.} Di kantin ada 5 jenis roti dan 3 jenis nasi. Andi hanya ingin membeli satu
makanan. Ada berapa pilihan?

\textit{Jawab:} Andi memilih roti \textbf{atau} nasi, jadi $5 + 3 = \textbf{8}$ pilihan.
\end{contoh}

\begin{giliranmu}
\begin{enumerate}[label=\alph*.]
    \item Di rak ada 7 novel dan 5 komik. Dina meminjam satu buku. Banyak pilihannya \kosong.
    \item Dari kota A ke kota B ada 2 jalur darat dan 3 jalur laut. Banyak cara pergi \kosong.
\end{enumerate}
\end{giliranmu}

\subsection*{Langkah 3: Teknik kotak}
Untuk menyusun bilangan, gambar kotak sebanyak digitnya, isi setiap kotak dengan banyak
pilihannya, lalu kalikan. Kalau angka tidak boleh berulang, pilihannya berkurang satu per satu.

\begin{contoh}
\textbf{Contoh.} Dari angka 1, 2, 3, 4, 5 dibuat bilangan 3 digit tanpa angka berulang.
Ada berapa bilangan?

\textit{Jawab:} Ratusan 5 pilihan, puluhan tinggal 4, satuan tinggal 3.
Jadi $5 \times 4 \times 3 = \textbf{60}$ bilangan.
\end{contoh}

\begin{giliranmu}
\begin{enumerate}[label=\alph*.]
    \item Bilangan 2 digit dari angka 1, 2, 3 (boleh berulang) ada \kosong.
    \item Bilangan 3 digit dari angka 1 sampai 6 tanpa angka berulang ada \kosong.
\end{enumerate}
\end{giliranmu}

\subsection*{Langkah 4: Kotak yang punya syarat}
Kotak yang punya syarat (misalnya harus genap) diisi \textbf{paling dulu}. Ingat juga, angka 0
tidak boleh berada di depan.

\begin{contoh}
\textbf{Contoh.} Dari angka 1, 2, 3, 4, 5 dibuat bilangan 3 digit ganjil tanpa angka berulang.
Ada berapa bilangan?

\textit{Jawab:} Satuan punya syarat, jadi diisi duluan: 1, 3, atau 5 (3 pilihan). Ratusan
tinggal 4 pilihan, puluhan tinggal 3. Jadi $3 \times 4 \times 3 = \textbf{36}$ bilangan.
\end{contoh}

\begin{giliranmu}
\begin{enumerate}[label=\alph*.]
    \item Bilangan genap 2 digit dari angka 1, 2, 3, 4 tanpa angka berulang ada \kosong.
    \item Bilangan 2 digit dari angka 0, 1, 2, 3 tanpa angka berulang ada \kosong.
\end{enumerate}
\end{giliranmu}

\begin{rangkuman}
\begin{enumerate}
    \item DAN (tahap demi tahap): kalikan.
    \item ATAU (pilih salah satu): jumlahkan.
    \item Teknik kotak: isi setiap kotak dengan banyak pilihannya, lalu kalikan.
    \item Kotak bersyarat diisi paling dulu. Angka 0 tidak boleh di depan.
\end{enumerate}
\end{rangkuman}

\textit{Kunci Giliranmu:} Langkah 1: 12, 24. \quad Langkah 2: 12, 5. \quad Langkah 3: 9, 120.
\quad Langkah 4: 6, 9.

\newpage
\subsection*{Latihan Soal}

\level{Level 1: Pemanasan}
\begin{enumerate}
  \item Rina punya 4 kaos dan 3 rok. Ada berapa pasangan kaos dan rok yang bisa ia pakai?
        \jawab{12}{$4 \times 3$.}
  \item Dari kota A ke kota B ada 3 jalan, dan dari kota B ke kota C ada 4 jalan. Ada berapa
        rute dari A ke C lewat B?
        \jawab{12}{$3 \times 4$.}
  \item \textbf{(Kode Loker)} Kode loker terdiri dari 1 huruf (A sampai E) lalu 1 angka
        (0 sampai 9), misalnya C7. Ada berapa kode yang mungkin?
        \jawab{50}{$5 \times 10$.}
  \item Tiga koin dilempar bersamaan. Ada berapa hasil yang mungkin?
        \jawab{8}{$2 \times 2 \times 2$. Cocok untuk latihan menggambar diagram pohon.}
  \item Di toko ada 4 jenis pensil dan 6 jenis pulpen. Dina hanya ingin membeli satu alat tulis.
        Ada berapa pilihan?
        \jawab{10}{$4 + 6$.}
\end{enumerate}

\level{Level 2: Eksplorasi}
\begin{enumerate}
  \setcounter{enumi}{5}
  \item Ada berapa bilangan 3 digit yang bisa dibuat dari angka 0, 1, 2, 3, 4 jika angkanya
        boleh berulang?
        \jawab{100}{Ratusan 4 pilihan (bukan 0), puluhan 5, satuan 5.}
  \item Ada berapa bilangan genap 3 digit yang bisa dibuat dari angka 1, 2, 3, 4, 5 tanpa
        angka berulang?
        \jawab{24}{Satuan 2 atau 4 (2 pilihan), lalu ratusan 4, puluhan 3.}
  \item \textbf{(PIN Ponsel)} PIN terdiri dari 4 digit (0--9) yang semuanya berbeda.
        Ada berapa PIN yang mungkin?
        \jawab{5040}{$10 \times 9 \times 8 \times 7$. Di sini 0 boleh di depan karena PIN bukan bilangan.}
  \item \textbf{(Kode Brankas)} Kode brankas terdiri dari 3 digit (0--9). Digit pertama ganjil,
        digit kedua genap, dan digit ketiga bilangan prima. Ada berapa kode yang mungkin?
        \jawab{100}{Ganjil 5 pilihan, genap 5 pilihan (0 termasuk genap), prima 4 pilihan (2, 3, 5, 7).}
  \item Ada berapa bilangan 3 digit yang semua digitnya ganjil?
        \jawab{125}{$5 \times 5 \times 5$.}
\end{enumerate}

\level{Level 3: Tantangan OSK}
\begin{enumerate}
  \setcounter{enumi}{10}
  \item Ada berapa bilangan genap 3 digit yang ketiga digitnya berbeda?
        \jawab{328}{Bagi dua kasus. Satuan 0: $9 \times 8 = 72$. Satuan 2, 4, 6, atau 8:
        $4 \times 8 \times 8 = 256$ (ratusan tidak boleh 0 dan tidak boleh sama dengan satuan).}
  \item Ada berapa bilangan 4 digit yang keempat digitnya berbeda dan habis dibagi 5?
        \jawab{952}{Satuan 0: $9 \times 8 \times 7 = 504$. Satuan 5: $8 \times 8 \times 7 = 448$.}
  \item \textbf{(Peta Harta Karun)} Dari A ke D bisa lewat B atau lewat C. Dari A ke B ada 2 jalan,
        B ke D ada 3 jalan, A ke C ada 3 jalan, dan C ke D ada 2 jalan. Ada juga 1 jalan rahasia
        langsung dari A ke D. Ada berapa rute dari A ke D?
        \jawab{13}{$2 \times 3 + 3 \times 2 + 1$.}
  \item \textbf{(Bendera Kelas)} Bendera terdiri dari 3 garis mendatar dan tersedia 4 warna.
        Dua garis yang bersebelahan tidak boleh sama warnanya. Ada berapa bendera yang bisa dibuat?
        \jawab{36}{Atas 4 pilihan, tengah 3, bawah 3 (asal beda dari tengah).}
  \item Ada berapa bilangan ganjil 3 digit yang digitnya berbeda dan kurang dari 500?
        \jawab{144}{Ratusan 1 atau 3: satuan 4 pilihan, puluhan 8, jadi $2 \times 4 \times 8 = 64$.
        Ratusan 2 atau 4: satuan 5 pilihan, puluhan 8, jadi $2 \times 5 \times 8 = 80$.}
\end{enumerate}

\level{Level 4: Setara OSP (Provinsi)}
\textit{Tulis jawaban lengkap dengan alasannya.}
\begin{enumerate}
  \setcounter{enumi}{15}
  \item Buktikan bahwa bilangan $2^a \times 3^b$ mempunyai tepat $(a+1)(b+1)$ faktor positif.
        \jawab{bukti}{Setiap faktor dari $2^a \times 3^b$ hanya boleh memuat faktor prima 2 dan 3,
        dengan pangkat yang tidak melebihi pangkat aslinya. Jadi setiap faktor berbentuk
        $2^x \times 3^y$ dengan $x$ dari $0$ sampai $a$ dan $y$ dari $0$ sampai $b$.\\
        Memilih $x$ ada $a + 1$ cara, \emph{dan} memilih $y$ ada $b + 1$ cara. Setiap pasangan
        $(x, y)$ menghasilkan faktor yang berbeda (karena faktorisasi prima tunggal). Dengan
        aturan perkalian, banyak faktornya $(a + 1)(b + 1)$. \hfill$\blacksquare$\\
        \emph{Ini adalah alasan rumus banyak faktor di Bab 1, Topik 5.}}
\end{enumerate}

\level{Level 5: Setara OSN (Nasional)}
\textit{Tulis jawaban lengkap dengan alasannya.}
\begin{enumerate}
  \setcounter{enumi}{16}
  \item Tentukan banyaknya bilangan 4 digit yang hasil kali digit-digitnya sama dengan 12.
        Jelaskan bahwa tidak ada yang terlewat.
        \jawab{36}{Tidak ada digit 0 (hasil kalinya akan 0). Digit yang membagi 12 hanya
        1, 2, 3, 4, 6. Cari semua kumpulan 4 digit yang hasil kalinya 12, diurutkan dari kecil:
        $\{1,1,2,6\}$, $\{1,1,3,4\}$, $\{1,2,2,3\}$. Tidak ada yang lain: kalau ada digit 6,
        tiga digit lainnya berhasil kali 2, yaitu $\{1,1,2\}$; kalau ada 4 (tanpa 6), sisanya
        berhasil kali 3, yaitu $\{1,1,3\}$; kalau tidak ada 4 dan 6, digitnya hanya 1, 2, 3, dan
        $12 = 2 \times 2 \times 3$ memberi $\{1,2,2,3\}$.\\
        Setiap kumpulan berisi 4 digit dengan satu digit kembar, jadi bisa disusun
        $\frac{4!}{2!} = 12$ cara. Totalnya $3 \times 12 = 36$.}
\end{enumerate}

\newpage
% ====================================================================
\section*{Topik 2: Menghitung dengan Cerdik}
% ====================================================================

\begin{catatanguru}
Banyak siswa akan menjawab ``dari halaman 5 sampai 10 ada 5 halaman''. Biarkan mereka menulis
nomornya satu per satu di Ayo Selidiki supaya menemukan sendiri aturan $b - a + 1$.
\end{catatanguru}

\begin{selidiki}
Rani membaca buku dari halaman 5 sampai halaman 10. Tuliskan semua nomor halaman yang ia baca:
\rule{6cm}{0.4pt}

Lengkapi tabel berikut (boleh dengan menulis nomornya).

\begin{center}
\renewcommand{\arraystretch}{1.6}
\begin{tabular}{|l|*{4}{>{\centering\arraybackslash}p{2cm}|}}
\hline
Halaman & 5 sampai 10 & 3 sampai 9 & 10 sampai 20 & 1 sampai 50\\ \hline
Banyak halaman & & & & \\ \hline
\end{tabular}
\end{center}

Sekarang jawab:
\begin{enumerate}
    \item Apakah banyak halaman dari 5 sampai 10 sama dengan $10 - 5$?
    \item Tebak aturan untuk menghitung banyak bilangan dari $a$ sampai $b$.
\end{enumerate}

Dari halaman 5 sampai 10 ada 6 halaman, bukan 5! Banyak bilangan dari $a$ sampai $b$ adalah
$b - a + 1$. Ini contoh menghitung dengan cerdik tanpa menulis semuanya.
\end{selidiki}

\subsection*{Langkah 1: Banyak bilangan dari $a$ sampai $b$}
Banyak bilangan dari $a$ sampai $b$ adalah $\mathbf{b - a + 1}$. Jangan lupa tambah 1!

\begin{contoh}
\textbf{Contoh.} Ada berapa bilangan dari 15 sampai 87?

\textit{Jawab:} $87 - 15 + 1 = \textbf{73}$.
\end{contoh}

\begin{giliranmu}
\begin{enumerate}[label=\alph*.]
    \item Banyak bilangan dari 100 sampai 250 adalah \kosong.
    \item Kursi bernomor 12 sampai 40. Banyak kursinya \kosong.
\end{enumerate}
\end{giliranmu}

\subsection*{Langkah 2: Banyak kelipatan}
Banyak kelipatan $k$ dari 1 sampai $n$ adalah hasil bagi $n \div k$ (sisanya dibuang).
Kalau tidak mulai dari 1, kurangi dengan kelipatan yang tidak ikut.

\begin{contoh}
\textbf{Contoh.} Ada berapa kelipatan 4 dari 30 sampai 100?

\textit{Jawab:} Kelipatan 4 dari 1 sampai 100 ada 25. Dari 1 sampai 29 ada 7. Jadi
$25 - 7 = \textbf{18}$.
\end{contoh}

\begin{giliranmu}
\begin{enumerate}[label=\alph*.]
    \item Banyak kelipatan 5 dari 1 sampai 100 adalah \kosong.
    \item Banyak kelipatan 3 dari 50 sampai 100 adalah \kosong.
\end{enumerate}
\end{giliranmu}

\subsection*{Langkah 3: Hitung yang tidak diminta}
Kadang lebih mudah menghitung yang \textbf{tidak} diminta, lalu mengurangkannya dari semua.
Cara ini cocok untuk soal yang ada kata ``paling sedikit satu''.

\begin{contoh}
\textbf{Contoh.} Ada berapa bilangan dari 1 sampai 1000 yang memuat angka 7?

\textit{Jawab:} Hitung yang \textbf{tidak} memuat angka 7. Tulis 1 sampai 999 sebagai 3 digit
(001 sampai 999). Setiap digit punya 9 pilihan (selain 7): $9 \times 9 \times 9 = 729$, tapi
000 ikut terhitung, jadi 728. Angka 1000 juga tidak memuat 7.
Jadi yang memuat angka 7 ada $1000 - 728 - 1 = \textbf{271}$.
\end{contoh}

\begin{giliranmu}
\begin{enumerate}[label=\alph*.]
    \item Banyak bilangan 2 digit yang memuat angka 3 adalah \kosong.
    \item Tiga koin dilempar. Banyak hasil yang memuat paling sedikit satu Gambar \kosong.
\end{enumerate}
\end{giliranmu}

\subsection*{Langkah 4: Dua kelompok yang beririsan}
Banyak (A atau B) = banyak A + banyak B $-$ banyak (A dan B).
Bagian irisan dikurangi karena terhitung dua kali.

\begin{contoh}
\textbf{Contoh.} Ada berapa bilangan dari 1 sampai 100 yang habis dibagi 3 atau 5?

\textit{Jawab:} Kelipatan 3 ada 33, kelipatan 5 ada 20, dan kelipatan 15 (habis dibagi keduanya)
ada 6. Jadi $33 + 20 - 6 = \textbf{47}$.

\begin{center}
\begin{tikzpicture}[scale=0.6]
\draw[thick] (0,0) circle (1.4);
\draw[thick, dashed] (1.8,0) circle (1.4);
\node at (-0.5,0) {27}; \node at (0.9,0) {6}; \node at (2.3,0) {14};
\node at (-0.6,1.8) {\small Kelipatan 3}; \node at (2.4,1.8) {\small Kelipatan 5};
\end{tikzpicture}
\end{center}
\end{contoh}

\begin{giliranmu}
\begin{enumerate}[label=\alph*.]
    \item Banyak bilangan dari 1 sampai 60 yang habis dibagi 2 atau 3 adalah \kosong.
    \item Dari 40 siswa, 25 ikut futsal, 20 ikut voli, dan 10 ikut keduanya. Yang tidak ikut
          keduanya \kosong{} siswa.
\end{enumerate}
\end{giliranmu}

\begin{rangkuman}
\begin{enumerate}
    \item Banyak bilangan dari $a$ sampai $b$ adalah $b - a + 1$.
    \item Banyak kelipatan $k$ dari 1 sampai $n$ adalah $n \div k$ (sisanya dibuang).
    \item ``Paling sedikit satu'': hitung semua dikurangi yang tidak diminta.
    \item A atau B = A + B $-$ (A dan B).
\end{enumerate}
\end{rangkuman}

\textit{Kunci Giliranmu:} Langkah 1: 151, 29. \quad Langkah 2: 20, 17. \quad Langkah 3: 18, 7.
\quad Langkah 4: 40, 5.

\newpage
\subsection*{Latihan Soal}

\level{Level 1: Pemanasan}
\begin{enumerate}
  \item Ada berapa bilangan bulat dari 23 sampai 77?
        \jawab{55}{$77 - 23 + 1$.}
  \item \textbf{(Membaca Novel)} Rani membaca novel dari halaman 12 sampai halaman 150.
        Berapa halaman yang ia baca?
        \jawab{139}{$150 - 12 + 1$.}
  \item Ada berapa kelipatan 7 dari 1 sampai 200?
        \jawab{28}{$200 \div 7 = 28$ sisa 4.}
  \item \textbf{(Tiang Bendera)} Di sepanjang jalan lurus 20 meter dipasang tiang bendera setiap
        2 meter, termasuk di kedua ujungnya. Ada berapa tiang?
        \jawab{11}{Tiang di posisi 0, 2, 4, \dots, 20.}
  \item Ada berapa bilangan genap dari 10 sampai 50?
        \jawab{21}{10, 12, \dots, 50 sama dengan $2 \times 5, \dots, 2 \times 25$. Jadi $25 - 5 + 1$.}
\end{enumerate}

\level{Level 2: Eksplorasi}
\begin{enumerate}
  \setcounter{enumi}{5}
  \item Ada berapa kelipatan 6 dari 100 sampai 300?
        \jawab{34}{Kelipatan 6 dari 1--300 ada 50, dari 1--99 ada 16.}
  \item Ada berapa bilangan dari 1 sampai 200 yang habis dibagi 2 atau 3?
        \jawab{133}{$100 + 66 - 33$.}
  \item Ada berapa bilangan dari 1 sampai 100 yang tidak habis dibagi 2 dan tidak habis dibagi 5?
        \jawab{40}{Yang habis dibagi 2 atau 5 ada $50 + 20 - 10 = 60$. Jadi $100 - 60$.}
  \item \textbf{(Uang Saku)} Ada berapa cara membayar Rp10.000 dengan lembaran Rp1.000, Rp2.000,
        dan Rp5.000? (Tidak harus memakai semua jenis uang.)
        \jawab{10}{Daftar berdasarkan banyak lembar Rp5.000. Dua lembar: 1 cara. Satu lembar:
        sisa Rp5.000 dibayar dengan 0, 1, atau 2 lembar Rp2.000, jadi 3 cara. Nol lembar:
        sisa Rp10.000 dibayar dengan 0 sampai 5 lembar Rp2.000, jadi 6 cara.}
  \item \textbf{(Survei Kelas)} Di kelas berisi 30 siswa, 18 suka Matematika, 15 suka IPA, dan
        7 suka keduanya. Berapa siswa yang tidak suka keduanya?
        \jawab{4}{$30 - (18 + 15 - 7)$.}
\end{enumerate}

\level{Level 3: Tantangan OSK}
\begin{enumerate}
  \setcounter{enumi}{10}
  \item Ada berapa bilangan 3 digit yang memuat paling sedikit satu angka 5?
        \jawab{252}{Semua bilangan 3 digit ada 900. Yang tanpa angka 5 ada $8 \times 9 \times 9 = 648$.}
  \item Empat koin dilempar bersamaan. Ada berapa hasil yang memuat paling sedikit satu sisi Gambar?
        \jawab{15}{$16 - 1$.}
  \item Ada berapa bilangan dari 1 sampai 1000 yang bukan bilangan kuadrat dan bukan bilangan
        pangkat tiga?
        \jawab{962}{Kuadrat ada 31 ($31^2 = 961$), pangkat tiga ada 10, yang keduanya sekaligus
        ada 3 (1, 64, 729). Jadi $1000 - (31 + 10 - 3)$.}
  \item Ada berapa bilangan 3 digit yang jumlah digitnya 5?
        \jawab{15}{Ratusan 5: 1 bilangan. Ratusan 4: 2. Ratusan 3: 3. Ratusan 2: 4. Ratusan 1: 5.}
  \item Ada berapa segitiga yang panjang sisinya bilangan bulat dan kelilingnya 12?
        \jawab{3}{Urutkan sisi $a \le b \le c$ dan pakai syarat $a + b > c$. Hasilnya
        (2, 5, 5), (3, 4, 5), dan (4, 4, 4).}
\end{enumerate}

\level{Level 4: Setara OSP (Provinsi)}
\textit{Tulis jawaban lengkap dengan alasannya.}
\begin{enumerate}
  \setcounter{enumi}{15}
  \item Buktikan bahwa banyak bilangan 3 digit yang jumlah digitnya 5 sama dengan banyak bilangan
        3 digit yang jumlah digitnya 23, \emph{tanpa} menghitung keduanya.
        \jawab{bukti}{Pasangkan bilangan $abc$ (jumlah digit 5) dengan bilangan
        $(10 - a)(9 - b)(9 - c)$. Jumlah digitnya $28 - (a + b + c) = 28 - 5 = 23$.\\
        Pasangan ini sah: karena $1 \le a \le 9$, maka $1 \le 10 - a \le 9$ (tetap digit depan yang
        bukan 0), dan $9 - b$, $9 - c$ tetap digit 0 sampai 9. Sebaliknya, setiap bilangan dengan
        jumlah digit 23 bisa dikembalikan dengan cara yang sama. Jadi setiap bilangan di kelompok
        pertama punya tepat satu pasangan di kelompok kedua, dan banyaknya sama. \hfill$\blacksquare$\\
        \emph{(Keduanya ada 15.) Ide memasangkan seperti ini disebut korespondensi satu-satu.}}
\end{enumerate}

\level{Level 5: Setara OSN (Nasional)}
\textit{Tulis jawaban lengkap dengan alasannya.}
\begin{enumerate}
  \setcounter{enumi}{16}
  \item Tentukan banyaknya bilangan dari 1 sampai 1000 yang tidak habis dibagi 2, 3, maupun 5.
        Jelaskan caramu.
        \jawab{266}{Hitung yang habis dibagi 2, 3, atau 5, dengan aturan tiga kelompok:
        tambahkan masing-masing, kurangi irisan dua-dua, lalu tambahkan irisan ketiganya.\\
        Habis dibagi 2: 500, dibagi 3: 333, dibagi 5: 200. Dibagi 6: 166, dibagi 10: 100,
        dibagi 15: 66. Dibagi 30: 33.\\
        Yang habis dibagi 2, 3, atau 5 ada $500 + 333 + 200 - 166 - 100 - 66 + 33 = 734$.
        Jadi yang tidak habis dibagi ketiganya ada $1000 - 734 = 266$.\\
        \emph{Irisan ketiganya ditambahkan lagi karena ikut terkurangi tiga kali saat
        mengurangi irisan dua-dua. Jelaskan dengan gambar Venn tiga lingkaran.}}
\end{enumerate}

\newpage
% ====================================================================
\section*{Topik 3: Permutasi (Menyusun)}
% ====================================================================

\begin{catatanguru}
Ayo Selidiki bisa diperagakan dengan 3 siswa yang berbaris di depan kelas. Untuk huruf kembar
(Langkah 3), tulis keenam susunan MAMA di papan supaya siswa melihat sendiri kenapa harus dibagi.
\end{catatanguru}

\begin{selidiki}
Tiga huruf A, B, C akan disusun berjajar. Tuliskan semua susunannya.

\begin{center}
\renewcommand{\arraystretch}{1.6}
\begin{tabular}{|*{6}{>{\centering\arraybackslash}p{1.5cm}|}}
\hline
ABC & ACB & & & & \\ \hline
\end{tabular}
\end{center}

Sekarang lengkapi tabel ini (coba tebak untuk 4 dan 5 huruf).

\begin{center}
\renewcommand{\arraystretch}{1.6}
\begin{tabular}{|l|*{5}{>{\centering\arraybackslash}p{1.2cm}|}}
\hline
Banyak huruf & 1 & 2 & 3 & 4 & 5\\ \hline
Banyak susunan & 1 & 2 & & & \\ \hline
\end{tabular}
\end{center}

Apa polanya? Coba lihat: $2 = 2 \times 1$, lalu $6 = 3 \times 2 \times 1$. Untuk 4 huruf,
posisi pertama ada 4 pilihan, posisi kedua 3, dan seterusnya, jadi $4 \times 3 \times 2 \times 1 = 24$.
\end{selidiki}

\subsection*{Langkah 1: Faktorial}
$n! = n \times (n-1) \times \dots \times 2 \times 1$, dibaca ``$n$ faktorial''.
Contoh: $4! = 24$, $5! = 120$, $6! = 720$. Menyusun $n$ benda berbeda ada $n!$ cara.

\begin{contoh}
\textbf{Contoh.} Lima sahabat berfoto berjajar. Ada berapa susunan?

\textit{Jawab:} $5! = 5 \times 4 \times 3 \times 2 \times 1 = \textbf{120}$ susunan.
\end{contoh}

\begin{giliranmu}
\begin{enumerate}[label=\alph*.]
    \item Banyak cara 6 orang antre di kasir adalah \kosong.
    \item Nilai $\dfrac{6!}{4!}$ adalah \kosong.
\end{enumerate}
\end{giliranmu}

\subsection*{Langkah 2: Memilih juara (urutan penting)}
Kalau hanya sebagian yang disusun, misalnya juara 1, 2, 3 dari 8 peserta, kalikan saja
$8 \times 7 \times 6$ (sebanyak posisi yang diisi).

\begin{contoh}
\textbf{Contoh.} Dari 8 peserta dipilih juara 1, 2, dan 3. Ada berapa kemungkinan?

\textit{Jawab:} $8 \times 7 \times 6 = \textbf{336}$.
\end{contoh}

\begin{giliranmu}
\begin{enumerate}[label=\alph*.]
    \item Ketua dan wakil ketua dipilih dari 10 siswa. Banyak caranya \kosong.
    \item PIN 3 digit berbeda dari angka 1 sampai 9 ada \kosong.
\end{enumerate}
\end{giliranmu}

\subsection*{Langkah 3: Huruf kembar}
Kalau ada huruf yang sama, bagi dengan faktorial banyaknya. Contoh MAMA: huruf M ada 2 dan
A ada 2, jadi $\dfrac{4!}{2! \times 2!} = 6$.

\begin{contoh}
\textbf{Contoh.} Ada berapa susunan huruf dari kata KATAK?

\textit{Jawab:} Ada 5 huruf: K dua kali, A dua kali, T sekali.
$\dfrac{5!}{2! \times 2!} = \dfrac{120}{4} = \textbf{30}$.
\end{contoh}

\begin{giliranmu}
\begin{enumerate}[label=\alph*.]
    \item Banyak susunan huruf MAMA adalah \kosong.
    \item Banyak susunan huruf ADA adalah \kosong.
\end{enumerate}
\end{giliranmu}

\subsection*{Langkah 4: Harus berdampingan dan duduk melingkar}
\textbf{Harus berdampingan:} ikat keduanya jadi satu, lalu ingat mereka bisa bertukar tempat.
\textbf{Duduk melingkar:} $n$ orang mengelilingi meja bundar ada $(n-1)!$ cara.

\begin{contoh}
\textbf{Contoh.} Ayah, Ibu, dan 3 anak berfoto berjajar. Ayah dan Ibu harus berdampingan.
Ada berapa susunan?

\textit{Jawab:} Anggap Ayah dan Ibu satu orang. Ada 4 ``orang'' yang disusun: $4! = 24$ cara.
Ayah dan Ibu bisa bertukar tempat (2 cara). Jadi $24 \times 2 = \textbf{48}$.
\end{contoh}

\begin{giliranmu}
\begin{enumerate}[label=\alph*.]
    \item 4 orang (termasuk A dan B) berjajar, A dan B berdampingan. Banyak susunannya \kosong.
    \item 4 orang duduk mengelilingi meja bundar. Banyak susunannya \kosong.
\end{enumerate}
\end{giliranmu}

\begin{rangkuman}
\begin{enumerate}
    \item Menyusun $n$ benda berbeda: $n!$ cara.
    \item Memilih juara 1, 2, 3 dari $n$: $n \times (n-1) \times (n-2)$.
    \item Huruf kembar: bagi dengan faktorial banyaknya huruf yang sama.
    \item Berdampingan: ikat jadi satu, kali 2. Melingkar: $(n-1)!$.
\end{enumerate}
\end{rangkuman}

\textit{Kunci Giliranmu:} Langkah 1: 720, 30. \quad Langkah 2: 90, 504. \quad Langkah 3: 6, 3.
\quad Langkah 4: 12, 6.

\newpage
\subsection*{Latihan Soal}

\level{Level 1: Pemanasan}
\begin{enumerate}
  \item Hitunglah $5!$ dan $\dfrac{7!}{5!}$.
        \jawab{120 dan 42}{$\frac{7!}{5!} = 7 \times 6$. Ajarkan untuk mencoret, bukan menghitung $5040 \div 120$.}
  \item Ada berapa cara menyusun 4 buku berbeda di rak?
        \jawab{24}{$4!$.}
  \item Dari 6 peserta lomba akan dipilih juara 1 dan juara 2. Ada berapa kemungkinan?
        \jawab{30}{$6 \times 5$.}
  \item Ada berapa susunan huruf dari kata BUKU?
        \jawab{12}{$\frac{4!}{2!}$, karena huruf U ada dua.}
  \item Dari 7 siswa akan dipilih ketua, sekretaris, dan bendahara. Ada berapa cara?
        \jawab{210}{$7 \times 6 \times 5$.}
\end{enumerate}

\level{Level 2: Eksplorasi}
\begin{enumerate}
  \setcounter{enumi}{5}
  \item Ada berapa susunan huruf dari kata MATEMATIKA?
        \jawab{151.200}{Ada 10 huruf: M dua kali, A tiga kali, T dua kali. Jadi $\frac{10!}{2! \times 3! \times 2!}$.}
  \item Dari susunan huruf KATAK, ada berapa yang diawali huruf K?
        \jawab{12}{Huruf K di depan, sisanya A, T, A, K disusun: $\frac{4!}{2!}$.}
  \item Enam siswa, termasuk Andi dan Budi, duduk berjajar. Andi dan Budi harus berdampingan.
        Ada berapa susunan?
        \jawab{240}{$5! \times 2$.}
  \item Sama seperti soal nomor 8, tetapi sekarang Andi dan Budi tidak boleh berdampingan.
        \jawab{480}{$6! - 240$. Ini memakai cara ``hitung yang tidak diminta'' dari Topik 2.}
  \item \textbf{(Antrean)} Tiga siswa putra dan tiga siswa putri antre berselang-seling
        (putra, putri, putra, \dots{} atau putri, putra, putri, \dots). Ada berapa susunan?
        \jawab{72}{Ada 2 pola. Setiap pola: putra $3!$, putri $3!$. Jadi $2 \times 6 \times 6$.}
\end{enumerate}

\level{Level 3: Tantangan OSK}
\begin{enumerate}
  \setcounter{enumi}{10}
  \item Lima orang duduk mengelilingi meja bundar. Ada berapa susunan?
        \jawab{24}{$4!$.}
  \item Dari angka 1, 2, 3, 4, 5 dibuat bilangan 5 digit tanpa angka berulang. Ada berapa
        bilangan yang lebih dari 30.000?
        \jawab{72}{Digit pertama 3, 4, atau 5, sisanya bebas: $3 \times 4!$.}
  \item \textbf{(Kamus Ajaib)} Semua susunan huruf A, B, C, D, E ditulis berurutan seperti di
        kamus: ABCDE, ABCED, dan seterusnya. CADEB ada di urutan ke berapa?
        \jawab{52}{Yang diawali A ada 24, diawali B ada 24. Yang diawali CA: CABDE, CABED, CADBE,
        CADEB. Jadi $48 + 4$.}
  \item Empat buku Matematika dan tiga buku IPA (semuanya berbeda) disusun di rak. Buku dengan
        pelajaran yang sama harus berkumpul. Ada berapa susunan?
        \jawab{288}{$2! \times 4! \times 3!$.}
  \item Lima siswa putra dan tiga siswa putri berbaris. Tidak boleh ada dua siswa putri yang
        berdampingan. Ada berapa susunan?
        \textit{(Petunjuk: susun putra dulu, lalu letakkan putri di sela-sela mereka.)}
        \jawab{14.400}{Putra disusun $5! = 120$ cara. Lima putra membuat 6 sela (termasuk kedua ujung).
        Tiga putri ditempatkan di 3 sela: $6 \times 5 \times 4 = 120$. Jadi $120 \times 120$.}
\end{enumerate}

\level{Level 4: Setara OSP (Provinsi)}
\textit{Tulis jawaban lengkap dengan alasannya.}
\begin{enumerate}
  \setcounter{enumi}{15}
  \item Buktikan bahwa banyak susunan huruf MATEMATIKA yang diawali huruf M sama dengan
        $\frac{1}{5}$ dari banyak seluruh susunannya.
        \jawab{bukti}{Seluruh susunan: $\frac{10!}{2! \times 3! \times 2!} = 151.200$.\\
        Susunan yang diawali M: huruf M pertama dikunci di depan, sisa 9 huruf (M, A, A, A, T, T,
        E, I, K) disusun: $\frac{9!}{3! \times 2!} = 30.240$.\\
        Perbandingannya $\frac{30.240}{151.200} = \frac{1}{5}$. \hfill$\blacksquare$\\
        \emph{Cara lain yang lebih cerdik: ada 10 huruf dan 2 di antaranya M. Kalau semua susunan
        ditulis, setiap posisi huruf sama ``adilnya'', jadi huruf pertama adalah M pada
        $\frac{2}{10} = \frac{1}{5}$ bagian dari semua susunan.}}
\end{enumerate}

\level{Level 5: Setara OSN (Nasional)}
\textit{Tulis jawaban lengkap dengan alasannya.}
\begin{enumerate}
  \setcounter{enumi}{16}
  \item Huruf A, A, B, B, C, C disusun berjajar. Tentukan banyak susunan yang tidak memiliki dua
        huruf sama yang berdampingan. Jelaskan caramu.
        \jawab{30}{Semua susunan: $\frac{6!}{2! \times 2! \times 2!} = 90$. Kurangi yang
        ``buruk''.\\
        Susunan dengan AA berdampingan: ikat AA, sisa 5 unit (AA, B, B, C, C):
        $\frac{5!}{2! \times 2!} = 30$. Sama untuk BB dan CC, jadi $3 \times 30 = 90$.\\
        Tetapi susunan dengan AA dan BB sama-sama berdampingan terhitung dua kali. Banyaknya:
        4 unit (AA, BB, C, C): $\frac{4!}{2!} = 12$. Ada 3 pasangan seperti ini, jadi $3 \times 12 = 36$.\\
        Susunan dengan AA, BB, CC semuanya berdampingan: $3! = 6$.\\
        Seperti Venn tiga lingkaran: yang buruk ada $90 - 36 + 6 = 60$. Jadi yang baik ada
        $90 - 60 = 30$.}
\end{enumerate}

\newpage
% ====================================================================
\section*{Topik 4: Kombinasi (Memilih)}
% ====================================================================

\begin{catatanguru}
Ayo Selidiki adalah inti topik ini. Pastikan siswa benar-benar mendaftar kedua daftar dan melihat
bahwa setiap pasangan perwakilan muncul dua kali di daftar ketua--wakil.
\end{catatanguru}

\begin{selidiki}
Ada 4 siswa: A, B, C, D.

\textbf{(1)} Pilih \textbf{ketua dan wakil}. Tulis semua kemungkinannya (ketua ditulis dulu):

AB, BA, AC, \rule{7cm}{0.4pt}

\textbf{(2)} Pilih \textbf{2 perwakilan} untuk lomba. Tulis semua kemungkinannya:

AB, AC, \rule{7cm}{0.4pt}

Sekarang jawab:
\begin{enumerate}
    \item Ada berapa kemungkinan di daftar (1)? Di daftar (2)?
    \item Kenapa daftar (2) tepat setengah dari daftar (1)?
\end{enumerate}

Pada daftar (1), AB dan BA berbeda (ketuanya beda). Pada daftar (2), AB dan BA sama saja.
Jadi setiap pasangan di daftar (2) muncul dua kali di daftar (1). Memilih tanpa memedulikan
urutan disebut \textbf{kombinasi}.
\end{selidiki}

\subsection*{Langkah 1: Rumus kombinasi}
Banyak cara memilih $r$ benda dari $n$ benda (urutan tidak penting) ditulis $\C{n}{r}$:
\[ \C{n}{r} = \frac{n \times (n-1) \times \dots \text{ (sebanyak $r$ faktor)}}{r!} \]

\begin{contoh}
\textbf{Contoh.} Dari 8 siswa dipilih 3 perwakilan. Ada berapa cara?

\textit{Jawab:} $\C{8}{3} = \dfrac{8 \times 7 \times 6}{3 \times 2 \times 1} = \textbf{56}$.
(Bandingkan dengan juara 1, 2, 3 yang ada 336 cara. Dibagi $3! = 6$ karena urutannya tidak penting.)
\end{contoh}

\begin{giliranmu}
\begin{enumerate}[label=\alph*.]
    \item Memilih 3 dari 6 siswa ada \kosong{} cara.
    \item Nilai $\C{10}{2}$ adalah \kosong.
\end{enumerate}
\end{giliranmu}

\subsection*{Langkah 2: Jabat tangan dan pertandingan}
Banyak pasangan dari $n$ orang adalah $\C{n}{2} = \dfrac{n(n-1)}{2}$. Rumus ini dipakai untuk
jabat tangan, pertandingan setengah kompetisi, dan garis yang menghubungkan dua titik.

\begin{contoh}
\textbf{Contoh.} Di pesta ada 10 orang. Setiap dua orang berjabat tangan sekali. Ada berapa
jabat tangan?

\textit{Jawab:} $\C{10}{2} = \dfrac{10 \times 9}{2} = \textbf{45}$.
\end{contoh}

\begin{giliranmu}
\begin{enumerate}[label=\alph*.]
    \item Ada 8 tim, setiap dua tim bertanding sekali. Banyak pertandingannya \kosong.
    \item Banyak diagonal segi-6 adalah \kosong. \quad (garis antar-titik sudut dikurangi sisi)
\end{enumerate}
\end{giliranmu}

\subsection*{Langkah 3: Memilih dari dua kelompok}
Kalau memilih dari dua kelompok sekaligus, pilih dari masing-masing kelompok lalu
\textbf{kalikan}. Untuk ``paling sedikit satu'', hitung semua dikurangi yang tidak diminta.

\begin{contoh}
\textbf{Contoh.} Dari 5 putra dan 4 putri dibentuk panitia berisi 2 putra dan 2 putri.
Ada berapa cara?

\textit{Jawab:} $\C{5}{2} \times \C{4}{2} = 10 \times 6 = \textbf{60}$.
\end{contoh}

\begin{giliranmu}
\begin{enumerate}[label=\alph*.]
    \item Dari 3 putra dan 4 putri dipilih 1 putra dan 2 putri. Banyak caranya \kosong.
    \item Dari 4 putra dan 3 putri dipilih 3 orang, paling sedikit 1 putri. Banyak caranya \kosong.
\end{enumerate}
\end{giliranmu}

\subsection*{Langkah 4: Jalan di grid}
Untuk berjalan dari pojok kiri bawah ke kanan atas dengan hanya melangkah ke kanan atau ke atas,
hitung total langkahnya, lalu pilih langkah mana yang ``atas''.

\begin{contoh}
\textbf{Contoh.} Ada berapa jalan dengan 4 langkah ke kanan dan 3 langkah ke atas?

\textit{Jawab:} Total 7 langkah, pilih 3 yang ``atas'': $\C{7}{3} = \textbf{35}$.
\end{contoh}

\begin{giliranmu}
\begin{enumerate}[label=\alph*.]
    \item Banyak jalan dengan 2 langkah ke kanan dan 2 langkah ke atas \kosong.
    \item Banyak jalan dengan 3 langkah ke kanan dan 2 langkah ke atas \kosong.
\end{enumerate}
\end{giliranmu}

\begin{rangkuman}
\begin{enumerate}
    \item Urutan penting (juara, jabatan): permutasi. Urutan tidak penting (tim, perwakilan): kombinasi.
    \item $\C{n}{r}$ = (kalikan $r$ bilangan dari $n$ turun) dibagi $r!$.
    \item Banyak pasangan dari $n$ orang: $\frac{n(n-1)}{2}$.
    \item Jalan di grid: pilih langkah mana yang ``atas''.
\end{enumerate}
\end{rangkuman}

\textit{Kunci Giliranmu:} Langkah 1: 20, 45. \quad Langkah 2: 28, 9. \quad Langkah 3: 18, 31.
\quad Langkah 4: 6, 10.

\newpage
\subsection*{Latihan Soal}

\level{Level 1: Pemanasan}
\begin{enumerate}
  \item Hitunglah $\C{6}{2}$ dan $\C{7}{5}$.
        \jawab{15 dan 21}{$\C{7}{5} = \C{7}{2}$, karena memilih 5 yang ikut sama dengan memilih 2 yang tidak ikut.}
  \item \textbf{(Es Krim)} Toko es krim punya 5 rasa. Kamu boleh memilih 2 rasa berbeda.
        Ada berapa pilihan?
        \jawab{10}{$\C{5}{2}$.}
  \item Dari 5 soal esai, siswa harus memilih 3 soal untuk dikerjakan. Ada berapa cara?
        \jawab{10}{$\C{5}{3}$.}
  \item \textbf{(Liga Futsal)} Ada 6 tim futsal. Setiap dua tim bertanding tepat satu kali.
        Ada berapa pertandingan?
        \jawab{15}{$\C{6}{2}$.}
  \item Ada 7 titik dan tidak ada 3 titik yang segaris. Ada berapa garis yang bisa dibuat
        melalui 2 titik?
        \jawab{21}{$\C{7}{2}$.}
\end{enumerate}

\level{Level 2: Eksplorasi}
\begin{enumerate}
  \setcounter{enumi}{5}
  \item Ada berapa diagonal pada segi-10?
        \jawab{35}{Semua garis antar-titik sudut $\C{10}{2} = 45$, dikurangi 10 sisi.}
  \item Ada 8 titik pada sebuah lingkaran. Ada berapa segitiga yang titik sudutnya dari
        titik-titik itu?
        \jawab{56}{$\C{8}{3}$.}
  \item Dari 5 siswa putra dan 4 siswa putri akan dipilih tim berisi 3 orang, dengan paling
        sedikit 1 putri. Ada berapa cara?
        \jawab{74}{Semua tim $\C{9}{3} = 84$, dikurangi tim tanpa putri $\C{5}{3} = 10$.}
  \item Ujian terdiri dari 10 soal dan siswa wajib mengerjakan 7 soal. Soal nomor 1, 2, dan 3
        wajib dikerjakan. Ada berapa pilihan soal?
        \jawab{35}{Sisa 4 soal dipilih dari 7 soal: $\C{7}{4}$.}
  \item \textbf{(Jalan ke Sekolah)} Sekolah Dina berada 5 blok ke timur dan 3 blok ke utara
        dari rumahnya. Dina hanya berjalan ke timur atau ke utara. Ada berapa rute?
        \jawab{56}{8 langkah, pilih 3 langkah ke utara: $\C{8}{3}$.}
\end{enumerate}

\level{Level 3: Tantangan OSK}
\begin{enumerate}
  \setcounter{enumi}{10}
  \item Dalam sebuah pertemuan terjadi 66 jabat tangan, dan setiap dua orang berjabat tangan
        tepat sekali. Berapa orang yang hadir?
        \jawab{12}{$\frac{n(n-1)}{2} = 66$, jadi $n(n-1) = 132 = 12 \times 11$.}
  \item Ada 5 titik pada garis $g$ dan 4 titik pada garis $h$, dan kedua garis itu sejajar.
        Ada berapa segitiga yang bisa dibentuk dari titik-titik tersebut?
        \jawab{70}{Dua titik dari $g$ dan satu dari $h$: $10 \times 4 = 40$. Dua titik dari $h$
        dan satu dari $g$: $6 \times 5 = 30$.}
  \item Ada berapa persegi panjang (termasuk persegi) pada papan $4 \times 4$ kotak?
        \jawab{100}{Persegi panjang ditentukan oleh 2 dari 5 garis mendatar dan 2 dari 5 garis
        tegak: $10 \times 10$.}
  \item Ada berapa bilangan 3 digit yang digitnya naik dari kiri ke kanan, seperti 147 atau 359?
        \jawab{84}{Pilih 3 digit berbeda dari 1--9. Setiap pilihan hanya punya satu susunan
        yang naik: $\C{9}{3}$.}
  \item \textbf{(Bagi Permen)} Enam permen yang sama dibagikan kepada 3 anak, dan setiap anak
        mendapat paling sedikit 1 permen. Ada berapa cara?
        \textit{(Petunjuk: jajarkan 6 permen, ada 5 celah di antaranya. Taruh 2 sekat.)}
        \jawab{10}{$\C{5}{2}$. Contoh: $\circ\circ|\circ|\circ\circ\circ$ artinya (2, 1, 3).}
\end{enumerate}

\level{Level 4: Setara OSP (Provinsi)}
\textit{Tulis jawaban lengkap dengan alasannya.}
\begin{enumerate}
  \setcounter{enumi}{15}
  \item Buktikan bahwa banyak diagonal pada segi-$n$ adalah $\dfrac{n(n-3)}{2}$.
        \jawab{bukti}{Cara 1: Setiap dua titik sudut membentuk garis, jadi ada $\C{n}{2} =
        \frac{n(n-1)}{2}$ garis. Dari garis-garis itu, $n$ di antaranya adalah sisi. Jadi banyak
        diagonal $\frac{n(n-1)}{2} - n = \frac{n(n-1) - 2n}{2} = \frac{n(n-3)}{2}$.\\
        Cara 2: Dari setiap titik sudut bisa ditarik diagonal ke semua titik sudut lain kecuali
        dirinya sendiri dan dua tetangganya, yaitu $n - 3$ diagonal. Untuk $n$ titik sudut:
        $n(n-3)$. Setiap diagonal terhitung dua kali (dari kedua ujungnya), jadi dibagi 2.
        \hfill$\blacksquare$}
\end{enumerate}

\level{Level 5: Setara OSN (Nasional)}
\textit{Tulis jawaban lengkap dengan alasannya.}
\begin{enumerate}
  \setcounter{enumi}{16}
  \item Pada sebuah lingkaran terdapat 8 titik. Setiap dua titik dihubungkan dengan tali busur,
        dan tidak ada tiga tali busur yang berpotongan di satu titik yang sama di dalam lingkaran.
        Tentukan banyak titik potong tali busur di dalam lingkaran. Jelaskan caramu.
        \jawab{70}{Kunci soal ini adalah melihat bahwa \emph{setiap titik potong berasal dari
        tepat 4 titik di lingkaran}. Dua tali busur berpotongan di dalam lingkaran hanya jika
        keempat ujungnya berbeda dan berselang-seling. Sebaliknya, dari 4 titik di lingkaran
        ada tepat satu cara memasangkan mereka menjadi dua tali busur yang berpotongan (yaitu
        kedua diagonal segi empatnya).\\
        Jadi setiap pilihan 4 titik menghasilkan tepat 1 titik potong, dan setiap titik potong
        berasal dari tepat 1 pilihan 4 titik. Banyak titik potong $= \C{8}{4} = 70$.}
\end{enumerate}

\newpage
% ====================================================================
\section*{Topik 5: Prinsip Sarang Merpati}
% ====================================================================

\begin{catatanguru}
Ayo Selidiki bisa diperagakan dengan pensil dan kotak sungguhan. Untuk soal ``paling sedikit harus
mengambil berapa'', selalu tanyakan ``Kalau kamu sangat sial, apa yang terambil?'' Untuk soal
pembuktian, biasakan siswa menyebut apa merpatinya dan apa sarangnya.
\end{catatanguru}

\begin{selidiki}
Masukkan \textbf{5 pensil} ke dalam \textbf{4 kotak} (boleh ada kotak yang kosong). Coba tiga
cara yang berbeda dan tulis isi setiap kotak.

\begin{center}
\renewcommand{\arraystretch}{1.6}
\begin{tabular}{|l|*{4}{>{\centering\arraybackslash}p{1.4cm}|}}
\hline
 & Kotak 1 & Kotak 2 & Kotak 3 & Kotak 4\\ \hline
Cara 1 & 2 & 1 & 1 & 1\\ \hline
Cara 2 & & & & \\ \hline
Cara 3 & & & & \\ \hline
\end{tabular}
\end{center}

Sekarang jawab:
\begin{enumerate}
    \item Di setiap cara, apakah selalu ada kotak yang berisi 2 pensil atau lebih?
    \item Bisakah kamu membuat semua kotak berisi paling banyak 1 pensil? Kenapa?
    \item Kalau ada 9 pensil dan 4 kotak, apakah pasti ada kotak yang berisi 3 pensil atau lebih?
\end{enumerate}

Kalau benda lebih banyak dari kotaknya, pasti ada kotak yang berisi lebih dari satu. Ide
sederhana ini disebut \textbf{Prinsip Sarang Merpati}: bendanya ``merpati'', kotaknya ``sarang''.
\end{selidiki}

\subsection*{Langkah 1: Merpati dan sarang}
Kalau merpatinya lebih banyak dari sarangnya, pasti ada sarang yang berisi paling sedikit
2 merpati. Kuncinya: tentukan dulu apa merpatinya dan apa sarangnya.

\begin{contoh}
\textbf{Contoh.} Buktikan bahwa di antara 13 orang, pasti ada 2 orang yang lahir di bulan yang sama.

\textit{Jawab:} Merpatinya 13 orang, sarangnya 12 bulan. Karena $13 > 12$, pasti ada bulan yang
``dihuni'' paling sedikit 2 orang.
\end{contoh}

\begin{giliranmu}
\begin{enumerate}[label=\alph*.]
    \item Paling sedikit berapa siswa supaya pasti ada 2 siswa yang angka terakhir nomor absennya
          sama? \kosong
    \item Ada 27 siswa di kelas. Apakah pasti ada 2 siswa yang huruf awal namanya sama?
          (Ada 26 huruf.) \kosong
\end{enumerate}
\end{giliranmu}

\subsection*{Langkah 2: Kasus paling sial}
Untuk soal ``paling sedikit harus mengambil berapa supaya \textbf{pasti}\dots'', bayangkan kamu
sedang sangat sial, lalu tambah 1.

\begin{contoh}
\textbf{Contoh.} Di laci gelap ada 10 kaos kaki merah dan 10 kaos kaki biru. Paling sedikit
berapa kaos kaki harus diambil supaya pasti dapat sepasang yang warnanya sama?

\textit{Jawab:} Paling sial: ambil 2, dapat merah dan biru. Kaos kaki ketiga pasti sama dengan
salah satunya. Jadi jawabannya $2 + 1 = \textbf{3}$.
\end{contoh}

\begin{giliranmu}
\begin{enumerate}[label=\alph*.]
    \item Kotak berisi bola dengan 4 warna. Supaya pasti dapat 2 bola sewarna, ambil \kosong{} bola.
    \item Kotak berisi 3 bola merah dan 5 bola biru. Supaya pasti dapat 1 bola biru, ambil
          \kosong{} bola.
\end{enumerate}
\end{giliranmu}

\subsection*{Langkah 3: Bentuk umum}
Kalau $N$ benda dimasukkan ke $k$ kotak, pasti ada kotak yang berisi paling sedikit
$N \div k$ benda, \textbf{dibulatkan ke atas}.

\begin{contoh}
\textbf{Contoh.} Di kelas ada 30 siswa. Paling sedikit berapa siswa yang pasti lahir di bulan
yang sama?

\textit{Jawab:} $30 \div 12 = 2{,}5$, dibulatkan ke atas menjadi \textbf{3}.
\end{contoh}

\begin{giliranmu}
\begin{enumerate}[label=\alph*.]
    \item 50 kelereng dimasukkan ke 7 kotak. Pasti ada kotak berisi paling sedikit \kosong.
    \item Ada 100 siswa. Pasti ada paling sedikit \kosong{} siswa yang lahir di bulan yang sama.
\end{enumerate}
\end{giliranmu}

\subsection*{Langkah 4: Sarang berupa sisa pembagian}
Di olimpiade, sarang yang paling sering dipakai adalah \textbf{sisa pembagian} (Bab 1, Topik 3).

\begin{contoh}
\textbf{Contoh.} Buktikan bahwa di antara 6 bilangan bulat sembarang, pasti ada dua bilangan
yang selisihnya habis dibagi 5.

\textit{Jawab:} Sisa pembagian oleh 5 hanya ada 5 kemungkinan (0, 1, 2, 3, 4). Ini sarangnya.
Bilangannya 6, ini merpatinya. Jadi pasti ada dua bilangan dengan sisa yang sama, dan selisih
keduanya habis dibagi 5.
\end{contoh}

\begin{giliranmu}
\begin{enumerate}[label=\alph*.]
    \item Paling sedikit berapa bilangan supaya pasti ada dua yang selisihnya habis dibagi 3? \kosong
    \item Paling sedikit berapa bilangan supaya pasti ada dua yang selisihnya habis dibagi 4? \kosong
\end{enumerate}
\end{giliranmu}

\begin{rangkuman}
\begin{enumerate}
    \item Merpati lebih banyak dari sarang: pasti ada sarang berisi 2 atau lebih.
    \item Tentukan dulu apa merpatinya dan apa sarangnya.
    \item ``Paling sedikit harus mengambil'': kasus paling sial, lalu tambah 1.
    \item $N$ benda di $k$ kotak: ada kotak berisi paling sedikit $N \div k$ (dibulatkan ke atas).
    \item Sarang favorit di olimpiade: sisa pembagian.
\end{enumerate}
\end{rangkuman}

\textit{Kunci Giliranmu:} Langkah 1: 11, ya. \quad Langkah 2: 5, 4. \quad Langkah 3: 8, 9.
\quad Langkah 4: 4, 5.

\newpage
\subsection*{Latihan Soal}

\level{Level 1: Pemanasan}
\begin{enumerate}
  \item Paling sedikit berapa orang supaya pasti ada 2 orang yang lahir di hari (Senin--Minggu)
        yang sama?
        \jawab{8}{7 hari sebagai sarang, jadi butuh $7 + 1$ orang.}
  \item Sebuah kotak berisi 5 bola merah, 5 putih, dan 5 hijau. Paling sedikit berapa bola harus
        diambil (tanpa melihat) supaya pasti ada 2 bola yang warnanya sama?
        \jawab{4}{Paling sial: merah, putih, hijau. Ditambah 1.}
  \item Di kelas berisi 30 siswa, nilai ulangan berupa bilangan bulat 0 sampai 10. Paling sedikit
        berapa siswa yang pasti mendapat nilai sama?
        \jawab{3}{Ada 11 kemungkinan nilai (0 sampai 10 itu 11 bilangan). $30 \div 11$ dibulatkan ke atas.}
  \item Paling sedikit berapa orang supaya pasti ada 2 orang yang ulang tahunnya di tanggal dan
        bulan yang sama? (29 Februari juga dihitung.)
        \jawab{367}{Ada 366 kemungkinan tanggal.}
  \item Jelaskan kenapa di antara 3 bilangan bulat sembarang, pasti ada 2 bilangan yang
        jumlahnya genap.
        \jawab{bukti}{Sarangnya ganjil dan genap (2 sarang), merpatinya 3 bilangan. Jadi ada dua
        bilangan yang sama ganjil-genapnya, dan jumlah keduanya genap.}
\end{enumerate}

\level{Level 2: Eksplorasi}
\begin{enumerate}
  \setcounter{enumi}{5}
  \item Sebuah kotak berisi 8 bola merah, 6 biru, dan 3 hijau. Paling sedikit berapa bola harus
        diambil supaya pasti ada 3 bola yang warnanya sama?
        \jawab{7}{Paling sial: 2 merah, 2 biru, 2 hijau. Ditambah 1.}
  \item Kotaknya sama seperti nomor 6. Paling sedikit berapa bola harus diambil supaya pasti
        mendapat 2 bola biru?
        \jawab{13}{Paling sial: 8 merah, 3 hijau, 1 biru. Ditambah 1. Banyak siswa menjawab 4.}
  \item \textbf{(Lemari Gelap)} Di lemari ada 12 kaos kaki hitam dan 12 kaos kaki putih yang
        tercampur. Paling sedikit berapa kaos kaki harus diambil supaya pasti mendapat sepasang
        kaos kaki hitam?
        \jawab{14}{Paling sial: 12 putih dan 1 hitam. Ditambah 1.}
  \item Jelaskan kenapa di antara 11 bilangan bulat sembarang, pasti ada 2 bilangan yang
        selisihnya habis dibagi 10.
        \jawab{bukti}{Sarangnya angka satuan (0--9), ada 10. Bilangannya 11, jadi ada dua dengan
        angka satuan sama, dan selisihnya berakhiran 0.}
  \item Sebuah sekolah memiliki 1000 siswa. Paling sedikit berapa siswa yang pasti ulang
        tahunnya di tanggal dan bulan yang sama?
        \jawab{3}{$1000 \div 366 = 2{,}7\ldots$, dibulatkan ke atas.}
\end{enumerate}

\level{Level 3: Tantangan OSK}
\begin{enumerate}
  \setcounter{enumi}{10}
  \item Dari bilangan 1, 2, 3, \dots, 10, paling sedikit berapa bilangan harus dipilih supaya
        pasti ada 2 bilangan yang jumlahnya 11?
        \jawab{6}{Sarangnya pasangan $\{1,10\}, \{2,9\}, \{3,8\}, \{4,7\}, \{5,6\}$. Memilih 1--5
        masih gagal, jadi butuh 6.}
  \item Lima titik diletakkan di dalam persegi berukuran $2 \times 2$. Buktikan bahwa pasti ada
        2 titik yang jaraknya tidak lebih dari $\sqrt{2}$.
        \jawab{bukti}{Bagi persegi menjadi 4 persegi kecil $1 \times 1$. Ada 5 titik, jadi dua titik
        berada di persegi kecil yang sama. Jarak terjauh di dalamnya adalah diagonal, yaitu $\sqrt{2}$.}
  \item Paling sedikit berapa bilangan harus dipilih dari 1 sampai 100 supaya pasti ada 2 bilangan
        yang selisihnya habis dibagi 7?
        \jawab{8}{Sarangnya sisa bagi 7 (ada 7). Memilih 1--7 masih gagal.}
  \item Buktikan bahwa di kelas berisi 40 siswa, pasti ada paling sedikit 4 siswa yang lahir di
        bulan yang sama.
        \jawab{bukti}{Kalau setiap bulan paling banyak 3 siswa, totalnya paling banyak 36. Padahal ada 40.}
  \item Di sebuah pesta yang dihadiri beberapa orang, ada yang saling berjabat tangan. Buktikan
        bahwa pasti ada 2 orang yang banyak jabat tangannya sama.
        \jawab{bukti}{Misal ada $n$ orang. Banyak jabat tangan seseorang antara 0 dan $n-1$.
        Tapi 0 dan $n-1$ tidak mungkin muncul bersamaan: kalau ada orang yang menjabat semua
        orang, tidak ada yang 0. Jadi hanya ada $n-1$ kemungkinan untuk $n$ orang.}
\end{enumerate}

\level{Level 4: Setara OSP (Provinsi)}
\textit{Tulis jawaban lengkap dengan alasannya.}
\begin{enumerate}
  \setcounter{enumi}{15}
  \item Buktikan bahwa dari 5 bilangan bulat sembarang, pasti ada 3 bilangan yang jumlahnya
        habis dibagi 3.
        \jawab{bukti}{Lihat sisa setiap bilangan jika dibagi 3: 0, 1, atau 2 (3 sarang).\\
        \emph{Kasus 1:} ketiga sisa (0, 1, 2) muncul. Ambil satu bilangan dari setiap sisa.
        Jumlah sisanya $0 + 1 + 2 = 3$, jadi jumlahnya habis dibagi 3.\\
        \emph{Kasus 2:} paling banyak 2 macam sisa yang muncul. Lima bilangan masuk ke paling
        banyak 2 sarang, jadi ada sarang yang berisi paling sedikit 3 bilangan ($5 \div 2$
        dibulatkan ke atas). Tiga bilangan dengan sisa sama $r$ jumlahnya bersisa $3r$, yang
        habis dibagi 3.\\
        Di kedua kasus selalu ada 3 bilangan yang jumlahnya habis dibagi 3. \hfill$\blacksquare$}
\end{enumerate}

\level{Level 5: Setara OSN (Nasional)}
\textit{Tulis jawaban lengkap dengan alasannya.}
\begin{enumerate}
  \setcounter{enumi}{16}
  \item Tentukan bilangan $n$ terkecil sehingga dari setiap $n$ bilangan yang dipilih dari
        $1, 2, \dots, 20$, pasti ada dua bilangan yang salah satunya membagi yang lain.
        Buktikan jawabanmu.
        \jawab{$n = 11$}{\emph{10 bilangan belum cukup:} pilih 11, 12, \dots, 20. Tidak ada yang
        membagi yang lain, karena kelipatan terkecil dari bilangan $\ge 11$ (selain dirinya)
        sudah $\ge 22$.\\
        \emph{11 bilangan pasti cukup:} setiap bilangan bisa ditulis sebagai
        $2^k \times (\text{bilangan ganjil})$, misalnya $12 = 2^2 \times 3$. Bagian ganjilnya
        adalah salah satu dari 1, 3, 5, \dots, 19, ada 10 kemungkinan (sarang). Dari 11 bilangan,
        pasti ada dua dengan bagian ganjil yang sama, misalnya $2^a \times m$ dan $2^b \times m$
        dengan $a < b$. Maka yang pertama membagi yang kedua. \hfill$\blacksquare$}
\end{enumerate}

\newpage
% ====================================================================
\section*{Topik 6: Peluang}
% ====================================================================

\begin{catatanguru}
Kalau ada dadu, mainkan teka-teki Rani dan Doni sekitar 30 kali dan catat hasilnya di papan.
Kesalahan yang sering muncul: menganggap jumlah dua dadu (2 sampai 12) sama-sama mungkin.
Tabel di Ayo Selidiki adalah cara terbaik untuk meluruskannya.
\end{catatanguru}

\begin{selidiki}
Dua dadu dilempar, lalu mata dadunya dijumlahkan. Lengkapi tabel jumlahnya.

\begin{center}
\renewcommand{\arraystretch}{1.4}
\begin{tabular}{|c|*{6}{>{\centering\arraybackslash}p{0.8cm}|}}
\hline
$+$ & 1 & 2 & 3 & 4 & 5 & 6\\ \hline
1 & 2 & 3 & 4 & 5 & 6 & 7\\ \hline
2 & 3 & & & & & \\ \hline
3 & & & & & & \\ \hline
4 & & & & & & \\ \hline
5 & & & & & & \\ \hline
6 & & & & & & \\ \hline
\end{tabular}
\end{center}

Sekarang jawab:
\begin{enumerate}
    \item Ada berapa kotak seluruhnya?
    \item Jumlah berapa yang paling sering muncul? Berapa kali jumlah 7 muncul? Jumlah 12?
    \item Rani menang kalau jumlahnya 7, Doni menang kalau jumlahnya 12. Apakah adil?
\end{enumerate}

Ada 36 hasil yang sama-sama mungkin. Jumlah 7 muncul 6 kali, jumlah 12 hanya sekali. Jadi
peluang Rani menang $\frac{6}{36}$, sedangkan Doni $\frac{1}{36}$. Tidak adil!
\end{selidiki}

\subsection*{Langkah 1: Rumus peluang}
\[ \text{Peluang} = \frac{\text{banyak hasil yang diharapkan}}{\text{banyak semua hasil yang mungkin}} \]
Peluang selalu bernilai antara 0 (mustahil) dan 1 (pasti terjadi).

\begin{contoh}
\textbf{Contoh.} Sebuah dadu dilempar. Berapa peluang muncul bilangan prima?

\textit{Jawab:} Bilangan primanya 2, 3, 5 (3 hasil) dari 6 hasil. Peluangnya
$\frac{3}{6} = \textbf{$\frac{1}{2}$}$.
\end{contoh}

\begin{giliranmu}
\begin{enumerate}[label=\alph*.]
    \item Kartu bernomor 1 sampai 10, diambil satu. Peluang terambil bilangan prima \kosong.
    \item Satu huruf diambil dari kata MATEMATIKA. Peluang terambil huruf A \kosong.
\end{enumerate}
\end{giliranmu}

\subsection*{Langkah 2: Dua dadu}
Untuk dua dadu, selalu gunakan tabel 36 hasil di Ayo Selidiki. Jangan menganggap jumlah 2 sampai
12 sama-sama mungkin.

\begin{contoh}
\textbf{Contoh.} Dua dadu dilempar. Berapa peluang jumlahnya 7?

\textit{Jawab:} Dari tabel, jumlah 7 muncul 6 kali. Peluangnya $\frac{6}{36} = \textbf{$\frac{1}{6}$}$.
\end{contoh}

\begin{giliranmu}
\begin{enumerate}[label=\alph*.]
    \item Peluang jumlah dua dadu sama dengan 4 adalah \kosong.
    \item Peluang jumlah dua dadu kurang dari 5 adalah \kosong.
\end{enumerate}
\end{giliranmu}

\subsection*{Langkah 3: Hitung yang tidak diminta}
Peluang tidak terjadi $= 1 -$ peluang terjadi. Cocok untuk soal ``paling sedikit satu''.

\begin{contoh}
\textbf{Contoh.} Tiga koin dilempar. Berapa peluang muncul paling sedikit satu Gambar?

\textit{Jawab:} Semua hasil $2 \times 2 \times 2 = 8$. Yang tidak ada Gambar sama sekali hanya
AAA (1 hasil). Peluangnya $1 - \frac{1}{8} = \textbf{$\frac{7}{8}$}$.
\end{contoh}

\begin{giliranmu}
\begin{enumerate}[label=\alph*.]
    \item Dua dadu dilempar. Peluang muncul paling sedikit satu angka 6 adalah \kosong.
    \item Empat koin dilempar. Peluang muncul paling sedikit satu Angka adalah \kosong.
\end{enumerate}
\end{giliranmu}

\subsection*{Langkah 4: Peluang dengan kombinasi}
Kalau mengambil beberapa benda sekaligus, hitung banyak caranya dengan kombinasi (Topik 4).

\begin{contoh}
\textbf{Contoh.} Dari kantong berisi 3 bola merah dan 2 bola putih diambil 2 bola sekaligus.
Berapa peluang keduanya merah?

\textit{Jawab:} Semua cara mengambil 2 dari 5 bola: $\C{5}{2} = 10$. Cara mengambil 2 merah:
$\C{3}{2} = 3$. Peluangnya \textbf{$\frac{3}{10}$}.
\end{contoh}

\begin{giliranmu}
\begin{enumerate}[label=\alph*.]
    \item Kantong berisi 4 bola merah dan 2 putih, diambil 2 sekaligus. Peluang keduanya putih \kosong.
    \item Dari 3 putra dan 2 putri dipilih 2 orang. Peluang terpilih 1 putra dan 1 putri \kosong.
\end{enumerate}
\end{giliranmu}

\begin{rangkuman}
\begin{enumerate}
    \item Peluang = banyak hasil yang diharapkan dibagi banyak semua hasil.
    \item Dua dadu: gunakan tabel 36 hasil.
    \item Peluang tidak terjadi $= 1 -$ peluang terjadi.
    \item Mengambil beberapa benda sekaligus: hitung dengan kombinasi.
\end{enumerate}
\end{rangkuman}

\textit{Kunci Giliranmu:} Langkah 1: $\frac{2}{5}$, $\frac{3}{10}$. \quad Langkah 2:
$\frac{1}{12}$, $\frac{1}{6}$. \quad Langkah 3: $\frac{11}{36}$, $\frac{15}{16}$. \quad
Langkah 4: $\frac{1}{15}$, $\frac{3}{5}$.

\newpage
\subsection*{Latihan Soal}

\level{Level 1: Pemanasan}
\begin{enumerate}
  \item Sebuah dadu dilempar. Berapa peluang muncul mata dadu genap?
        \jawab{$\frac{1}{2}$}{$\frac{3}{6}$.}
  \item Dalam kantong ada 3 kelereng merah, 5 biru, dan 2 hijau. Diambil 1 kelereng. Berapa peluang
        terambil kelereng biru?
        \jawab{$\frac{1}{2}$}{$\frac{5}{10}$.}
  \item Kartu bernomor 1 sampai 20 dikocok, lalu diambil satu. Berapa peluang terambil kartu
        kelipatan 3?
        \jawab{$\frac{3}{10}$}{Ada 6 kartu: 3, 6, 9, 12, 15, 18.}
  \item Dua koin dilempar. Berapa peluang keduanya menunjukkan sisi Gambar?
        \jawab{$\frac{1}{4}$}{Hasilnya AA, AG, GA, GG.}
  \item Satu huruf diambil secara acak dari kata OLIMPIADE. Berapa peluang terambil huruf vokal?
        \jawab{$\frac{5}{9}$}{Vokalnya O, I, I, A, E.}
\end{enumerate}

\level{Level 2: Eksplorasi}
\begin{enumerate}
  \setcounter{enumi}{5}
  \item Dua dadu dilempar. Berapa peluang jumlah mata dadunya 10?
        \jawab{$\frac{1}{12}$}{(4,6), (5,5), (6,4): $\frac{3}{36}$.}
  \item Dua dadu dilempar. Berapa peluang kedua mata dadu sama?
        \jawab{$\frac{1}{6}$}{$\frac{6}{36}$.}
  \item Dua dadu dilempar. Berapa peluang hasil kali kedua mata dadu genap?
        \jawab{$\frac{3}{4}$}{Hasil kali ganjil hanya jika keduanya ganjil: $3 \times 3 = 9$ hasil.
        Jadi $1 - \frac{9}{36}$. Ini memakai paritas dari Bab 1.}
  \item Dari kantong berisi 4 bola merah dan 6 bola biru diambil 3 bola sekaligus. Berapa peluang
        paling sedikit satu bola merah?
        \jawab{$\frac{5}{6}$}{Peluang semua biru $\frac{\C{6}{3}}{\C{10}{3}} = \frac{20}{120} = \frac{1}{6}$.}
  \item Sebuah bilangan 2 digit dipilih secara acak. Berapa peluang bilangan itu habis dibagi 7?
        \jawab{$\frac{13}{90}$}{Ada 90 bilangan 2 digit, dan kelipatan 7 dari 14 sampai 98 ada 13.}
\end{enumerate}

\level{Level 3: Tantangan OSK}
\begin{enumerate}
  \setcounter{enumi}{10}
  \item Tiga dadu dilempar. Berapa peluang jumlah mata dadunya 5?
        \jawab{$\frac{1}{36}$}{Semua hasil 216. Jumlah 5 didapat dari susunan (1,1,3) sebanyak 3 dan
        (1,2,2) sebanyak 3. Jadi $\frac{6}{216}$.}
  \item Huruf-huruf kata KATAK disusun secara acak. Berapa peluang kedua huruf K berdampingan?
        \jawab{$\frac{2}{5}$}{Semua susunan 30. Kalau KK diikat: $\frac{4!}{2!} = 12$.}
  \item Empat orang, termasuk Andi dan Budi, duduk berjajar secara acak. Berapa peluang Andi dan
        Budi duduk berdampingan?
        \jawab{$\frac{1}{2}$}{$\frac{3! \times 2}{4!} = \frac{12}{24}$.}
  \item Dua dadu dilempar. Berapa peluang angka satuan dari hasil kali kedua mata dadu adalah 6?
        \jawab{$\frac{1}{6}$}{$1 \times 6$, $6 \times 1$, $2 \times 3$, $3 \times 2$, $4 \times 4$,
        $6 \times 6$: ada 6 hasil.}
  \item \textbf{(Hari Lahir)} Ada 3 orang. Berapa peluang paling sedikit 2 orang di antaranya lahir
        di hari (Senin--Minggu) yang sama?
        \jawab{$\frac{19}{49}$}{Semua kemungkinan $7^3 = 343$. Ketiganya beda hari:
        $7 \times 6 \times 5 = 210$. Jadi $1 - \frac{210}{343} = \frac{133}{343} = \frac{19}{49}$.}
\end{enumerate}

\level{Level 4: Setara OSP (Provinsi)}
\textit{Tulis jawaban lengkap dengan alasannya.}
\begin{enumerate}
  \setcounter{enumi}{15}
  \item Sebuah koin dilempar $n$ kali ($n \ge 1$). Buktikan bahwa peluang banyaknya Gambar genap
        (termasuk 0) sama dengan $\frac{1}{2}$.
        \jawab{bukti}{Pasangkan setiap hasil dengan hasil yang sama persis kecuali lemparan
        \emph{terakhir} dibalik (A menjadi G atau sebaliknya). Contoh untuk $n = 3$: AGA dipasangkan
        dengan AGG.\\
        Dalam setiap pasangan, banyak Gambar berbeda tepat 1, jadi yang satu genap dan yang lain
        ganjil. Semua $2^n$ hasil terbagi habis menjadi pasangan-pasangan seperti ini. Jadi hasil
        dengan banyak Gambar genap tepat setengah dari semua hasil, dan peluangnya $\frac{1}{2}$.
        \hfill$\blacksquare$}
\end{enumerate}

\level{Level 5: Setara OSN (Nasional)}
\textit{Tulis jawaban lengkap dengan alasannya.}
\begin{enumerate}
  \setcounter{enumi}{16}
  \item Tiga dadu dilempar. Tentukan peluang bahwa mata dadu \textbf{terbesar} adalah 4.
        Jelaskan caramu.
        \jawab{$\frac{37}{216}$}{Semua hasil $6^3 = 216$.\\
        ``Terbesar adalah 4'' artinya semua dadu paling besar 4, \emph{dan} ada paling sedikit satu
        dadu yang tepat 4.\\
        Semua dadu paling besar 4: setiap dadu punya 4 pilihan (1--4), jadi $4^3 = 64$ hasil.
        Dari hasil ini, yang tidak memuat 4 sama sekali (semua dadu paling besar 3) ada $3^3 = 27$.
        Jadi yang terbesarnya tepat 4 ada $64 - 27 = 37$.\\
        Peluangnya $\frac{37}{216}$.}
\end{enumerate}

\newpage
% ====================================================================
\section*{Soal Campuran}
% ====================================================================

Soal-soal di bawah ini sengaja tidak diberi keterangan topik, sama seperti soal olimpiade
sungguhan. Sebelum menghitung, tanya dulu ke dirimu sendiri: urutannya penting atau tidak?
Ada kata ``paling sedikit'' atau tidak? Waktu mengerjakan 60 menit.

\begin{enumerate}
  \item Ada berapa bilangan 4 digit yang semua digitnya genap?
        \jawab{500}{$4 \times 5 \times 5 \times 5$ (ribuan tidak boleh 0).}
  \item Ada 7 titik pada sebuah lingkaran. Ada berapa segi empat yang titik sudutnya dari
        titik-titik itu?
        \jawab{35}{$\C{7}{4}$.}
  \item Ada berapa susunan huruf dari kata ANANDA?
        \jawab{60}{$\frac{6!}{3! \times 2!}$.}
  \item Paling sedikit berapa siswa harus dikumpulkan supaya pasti ada 3 siswa yang huruf awal
        namanya sama? (Ada 26 huruf.)
        \jawab{53}{Paling sial: setiap huruf dipakai 2 siswa, yaitu 52. Ditambah 1.}
  \item Ada berapa bilangan dari 1 sampai 500 yang habis dibagi 4 atau 6?
        \jawab{167}{$125 + 83 - 41$. Irisannya kelipatan 12, bukan 24.}
  \item Dua dadu dilempar. Berapa peluang jumlah mata dadunya bilangan prima?
        \jawab{$\frac{5}{12}$}{Jumlah 2, 3, 5, 7, 11 muncul $1 + 2 + 4 + 6 + 2 = 15$ kali.}
  \item Delapan siswa dibagi menjadi dua tim berisi 4 orang. Tim tidak diberi nama.
        Ada berapa cara?
        \jawab{35}{$\C{8}{4} = 70$, dibagi 2.}
  \item \textbf{(Tangga Sekolah)} Doni menaiki 6 anak tangga. Setiap langkah ia naik 1 atau 2 anak
        tangga. Ada berapa cara?
        \jawab{13}{Mulai dari tangga kecil: 1, 2, 3, 5, 8, 13. Setiap nilai adalah jumlah dua
        nilai sebelumnya.}
  \item Ada berapa bilangan 3 digit yang digitnya berbeda dan habis dibagi 5?
        \jawab{136}{Satuan 0: 72. Satuan 5: 64.}
  \item \textbf{(Turnamen Gugur)} Ada 64 tim dalam turnamen sistem gugur. Tim yang kalah langsung
        tersingkir. Ada berapa pertandingan sampai juaranya ditemukan?
        \jawab{63}{Setiap pertandingan menyingkirkan 1 tim, dan 63 tim harus tersingkir.}
\end{enumerate}

\end{document}
